\documentclass[12pt,a4paper]{article}

\usepackage[utf8]{inputenc}
\usepackage[T2A]{fontenc}
\usepackage[russian]{babel}
\usepackage{lmodern}
\usepackage{mathtext}
\usepackage{microtype}
\usepackage{amsmath}
\usepackage{amssymb}
\usepackage{graphicx}
\usepackage{booktabs}
\usepackage{xcolor}
\usepackage{hyperref}
\usepackage{tikz}
\usepackage{float}
\usepackage{enumitem}
\usepackage{listings}
\usepackage{float}
\usepackage{placeins}
\usepackage[left=2cm,right=2cm,top=2cm]{geometry}

\usetikzlibrary{positioning}
\usetikzlibrary{arrows}

% Define JSON language for listings
\lstdefinelanguage{json}{
    basicstyle=\ttfamily\small,
    numbers=left,
    numberstyle=\tiny,
    stepnumber=1,
    numbersep=8pt,
    showstringspaces=false,
    breaklines=true,
    frame=lines,
    commentstyle=\color{gray}\upshape,
    stringstyle=\color{purple},
    keywordstyle=\color{blue},
    identifierstyle=\color{black},
    keywords={true, false, null},
    morestring=[b]",
    morecomment=[s]{/*}{*/},
    morecomment=[l]//,
    morecomment=[s]{/**}{*/},
    sensitive=true
}

\usepackage{titlesec}
\titleformat{\section}
{\normalfont\Large\bfseries}{\thesection}{1em}{}
\titleformat{\subsection}
{\normalfont\large\bfseries}{\thesubsection}{1em}{}

\lstset{
  basicstyle=\ttfamily\small,
  columns=flexible,
  breaklines=true,
  commentstyle=\color{gray},
  keywordstyle=\color{blue},
  stringstyle=\color{purple},
  numbers=left,
  numberstyle=\tiny\color{gray},
  stepnumber=1,
  numbersep=5pt,
  showstringspaces=false,
  captionpos=b
}

\hypersetup{
  colorlinks=true,
  linkcolor=blue,
  filecolor=magenta,
  urlcolor=cyan,
}

\begin{document}

\begin{titlepage}
  \begin{center}
    \vspace*{5cm}
    {\Huge \textbf{Индексатор Содержания Видеолекций}}\\
    \vspace{1.5cm}
    {\LARGE Техническое Задание \& План Реализации}\\
    \vspace{3cm}
    {\large \today}\\
    \vfill
  \end{center}

  \begin{flushright}
    \textit{Подготовил} \\
    студент 3 курса \\
    группы Б22-221 \\
    Сидюк Д.А.
  \end{flushright}
  \vspace{1cm}
\end{titlepage}

\tableofcontents
\newpage

\section{Краткое Изложение}

Индексатор содержания видеолекций — это система, предназначенная для обработки множества видеолекций с YouTube по различным академическим дисциплинам, создающая единый поисковый индекс концепций. Система позволяет студентам быстро находить соответствующий образовательный контент в большой базе видео, различая мимолётные упоминания и содержательное освещение концепций.

Основным техническим новшеством системы является применение алгоритмов линейного поиска множественной наибольшей общей подпоследовательности (MLCS) в сочетании с методами обработки естественного языка для идентификации, анализа и индексации академических концепций в транскриптах видео. Система будет интегрирована с параллельно разрабатываемой студенческой платформой, построенной на базе Java Spring с TypeScript/React фронтендом.

\section{Обзор Системы}

\subsection{Цели}

\begin{itemize}
    \item Создание единого индекса академических концепций по нескольким видеолекциям
    \item Различение между мимолётными упоминаниями и содержательными объяснениями концепций
    \item Обеспечение эффективного поиска конкретных концепций с точными временными метками видео
    \item Поддержка автоматизированной обработки нового видеоконтента с минимальным ручным вмешательством
    \item Обеспечение бесшовной интеграции с существующей образовательной платформой
\end{itemize}

\subsection{Обзор Архитектуры}

Система следует модульной архитектуре со специализированными компонентами для каждой фазы конвейера обработки:

\begin{figure}[H]
\centering
\begin{tikzpicture}[
    node distance=0.6cm,
    auto,
    block/.style={rectangle, draw, fill=blue!10, text width=3.1cm, text centered, rounded corners, minimum height=0.8cm},
    line/.style={draw, -stealth, thick}
]
    % Nodes
    \node [block] (data) {Сбор Данных};
    \node [block, right=of data] (preproc) {Предобработка Текста};
    \node [block, right=of preproc] (concept) {Извлечение Концепций};
    \node [block, right=of concept] (score) {Оценка Релевантности};
    \node [block, below=of data, yshift=-0.2cm] (index) {Индексация};
    \node [block, right=of index] (search) {Поиск \& Извлечение};
    \node [block, right=of search] (integrate) {Интеграционный Слой};

    % Connections
    \path [line] (data) -- (preproc);
    \path [line] (preproc) -- (concept);
    \path [line] (concept) -- (score);
    \path [line] (score) |- ([yshift=-0.3cm]score.south) -| (index);
    \path [line] (index) -- (search);
    \path [line] (search) -- (integrate);
\end{tikzpicture}
\caption{Обзор архитектуры системы}
\end{figure}

\newpage
\subsection{Ключевые Технологии}

\begin{itemize}
    \item \textbf{Языки программирования}: C++ (основные алгоритмы), Python (обработка данных, NLP)
    \item \textbf{Внешние API}: YouTube Data API v3
    \item \textbf{Библиотеки NLP}: spaCy, NLTK
    \item \textbf{База данных}: PostgreSQL (прод), SQLite (разработка)
    \item \textbf{Интеграция}: REST API, обмен данными в формате JSON
\end{itemize}

\newpage
\section{Технические Спецификации}

\subsection{Компонент Сбора Данных}

\subsubsection{Функциональность}
Компонент сбора данных отвечает за получение метаданных видео и транскриптов с YouTube.

\subsubsection{Технические Детали}
\begin{itemize}
    \item \textbf{Интеграция с YouTube API}
    \begin{itemize}
        \item Реализация Python-клиента для YouTube Data API v3
        \item Поддержка как ручной, так и пакетной отправки видео
        \item Обработка ограничений квоты API с дросселированием запросов
    \end{itemize}

    \item \textbf{Извлечение Транскрипта}
    \begin{itemize}
        \item Извлечение сегментов транскрипта с синхронизацией по времени
        \item Поддержка нескольких языковых опций с приоритизацией
        \item Обработка автоматически созданных субтитров с фильтрацией по уровню достоверности
    \end{itemize}

    \item \textbf{Обработка Метаданных}
    \begin{itemize}
        \item Извлечение заголовка видео, описания, тегов, продолжительности
        \item Идентификация специфических метаданных лекции при их наличии
        \item Хранение метаданных в структурированном формате, связанном с ID видео
    \end{itemize}
\end{itemize}

\subsection{Компонент Предобработки Текста}

\subsubsection{Функциональность}
Компонент предобработки текста нормализует и улучшает текст транскрипта для дальнейшего анализа.

\subsubsection{Технические Детали}
\begin{itemize}
    \item \textbf{Нормализация Текста}
    \begin{itemize}
        \item Очистка и нормализация текста транскрипта
        \item Сегментация текста на предложения
        \item Сохранение синхронизации с временными метками
    \end{itemize}

    \item \textbf{Лингвистический Анализ}
    \begin{itemize}
        \item Выполнение частеречной разметки
        \item Применение лемматизации с сохранением исходного текста
        \item Идентификация академических дискурсивных маркеров
    \end{itemize}

    \item \textbf{Структурный Анализ}
    \begin{itemize}
        \item Обнаружение структуры лекции (введение, основное содержание, заключение)
        \item Определение границ разделов на основе дискурсивных маркеров
        \item Создание структурного слоя метаданных
    \end{itemize}
\end{itemize}

\subsubsection{Конвейер Обработки}
\begin{figure}[H]
\centering
\begin{tikzpicture}[
    node distance=0.3cm,
    auto,
    block/.style={rectangle, draw, fill=blue!10, text width=4.2cm, text centered, rounded corners, minimum height=0.7cm},
    line/.style={draw, -stealth, thick}
]
    % First row
    \node [block] (input) {Входные данные: Сегменты транскрипта};
    \node [block, right=of input] (clean) {Очистка и нормализация текста};
    \node [block, right=of clean] (segment) {Сегментация предложений};

    % Second row
    \node [block, below=of input, yshift=-0.2cm] (pos) {Частеречная разметка};
    \node [block, right=of pos] (marker) {Идентификация дискурсивных маркеров};
    \node [block, right=of marker] (struct) {Структурный анализ};

    % Output
    \node [block, below=of pos, xshift=2.5cm, yshift=-0.2cm, text width=9cm] (output) {Выходные данные: Предобработанный транскрипт с аннотациями};

    % Connections
    \path [line] (input) -- (clean);
    \path [line] (clean) -- (segment);
    \path [line] (segment) |- ([yshift=-0.25cm]segment.south) -| (pos);
    \path [line] (pos) -- (marker);
    \path [line] (marker) -- (struct);
    \path [line] (struct) |- ([yshift=-0.45cm]struct.south) -| (output);
\end{tikzpicture}
\caption{Конвейер предобработки текста}
\end{figure}

\subsection{Компонент Извлечения Концепций}

\subsubsection{Функциональность}
Компонент извлечения концепций идентифицирует академические концепции в предобработанных транскриптах.

\subsubsection{Технические Детали}
\begin{itemize}
    \item \textbf{Идентификация Концепций по Предметным Областям}
    \begin{itemize}
        \item Реализация пользовательского распознавания именованных сущностей для академических областей
        \item Применение TF-IDF и TextRank для извлечения терминологии
        \item Обнаружение многословных академических концепций
    \end{itemize}

    \item \textbf{Алгоритм Linear MLCS}
    \begin{itemize}
        \item Реализация оптимизированной версии линейного MLCS на C++
        \item Определение общих последовательностей концепций в разных лекциях
        \item Обнаружение зависимостей концепций и предварительных условий
    \end{itemize}

    \item \textbf{Контекстный Анализ}
    \begin{itemize}
        \item Извлечение контекста, окружающего каждую концепцию
        \item Анализ лингвистических паттернов, указывающих на объяснение
        \item Измерение стабильности концепции в контексте
    \end{itemize}
\end{itemize}

\subsection{Компонент Оценки Релевантности}

\subsubsection{Функциональность}
Компонент оценки релевантности определяет глубину освещения для каждого вхождения концепции.

\subsubsection{Технические Детали}
\begin{itemize}
    \item \textbf{Анализ Глубины Освещения}
    \begin{itemize}
        \item Реализация анализа скользящего окна
        \item Расчет плотности концепции в локальном контексте
        \item Измерение устойчивого внимания к концепциям
    \end{itemize}

    \item \textbf{Оценка Академической Значимости}
    \begin{itemize}
        \item Анализ паттернов акцентирования внимания преподавателем
        \item Обнаружение педагогических структур
        \item Измерение центральности в структуре лекции
    \end{itemize}

    \item \textbf{Расчет Композитной Релевантности}
    \begin{itemize}
        \item Объединение нескольких факторов оценки
        \item Применение нормализации по всему корпусу видео
        \item Классификация вхождений по категориям релевантности
    \end{itemize}
\end{itemize}

\subsubsection{Категории Релевантности}

\begin{table}[h]
\centering
\begin{tabular}{@{}llp{8cm}@{}}
\toprule
\textbf{Категория} & \textbf{Диапазон Оценки} & \textbf{Описание} \\
\midrule
Содержательное Освещение & 70-100 & Подробное объяснение с определениями, примерами и применениями \\
Умеренное Освещение & 40-69 & Предоставляется некоторое объяснение, но не исчерпывающее \\
Краткое Упоминание & 10-39 & Концепция упоминается с минимальным контекстом \\
Мимолетное Упоминание & 0-9 & Концепция упоминается без объяснения \\
\bottomrule
\end{tabular}
\caption{Категории оценки релевантности}
\end{table}

\subsection{Компонент Индексации}

\subsubsection{Функциональность}
Компонент индексации создает эффективный поисковый индекс концепций по всем видео.

\subsubsection{Технические Детали}
\begin{itemize}
    \item \textbf{Построение Инвертированного Индекса}
    \begin{itemize}
        \item Реализация оптимизированного инвертированного индекса на C++
        \item Создание списков вхождений с временными метками и релевантностью
        \item Применение сжатия для эффективного хранения
    \end{itemize}

    \item \textbf{Сохранение Индекса}
    \begin{itemize}
        \item Реализация сериализации/десериализации
        \item Поддержка дифференциальных обновлений для нового контента
        \item Обеспечение потокобезопасного параллельного доступа
    \end{itemize}

    \item \textbf{Интеграция Графа Знаний}
    \begin{itemize}
        \item Преобразование взаимосвязей концепций в структуру графа
        \item Реализация возможностей навигации по концепциям
        \item Поддержка семантических запросов отношений
    \end{itemize}
\end{itemize}

\subsubsection{Структура Индекса}
\begin{lstlisting}[language=json, caption=Структура инвертированного индекса]
ConceptIndex {
  "concept_id": {
    "canonical_name": "weierstrass theorem",
    "variants": ["weierstrass theorem", "extreme value theorem"],
    "concept_type": "theorem",
    "domain": "mathematical_analysis",
    "occurrence_count": 42,
    "postings": [
      {
        "video_id": "abc123",
        "relevance_score": 89,
        "coverage_type": "substantive",
        "timestamp_start": 1250.5,
        "timestamp_end": 1380.2,
        "segment_id": "abc123_seg42"
      },
      // More occurrences...
    ]
  },
  // More concepts...
}
\end{lstlisting}

\subsection{Компонент Поиска и Извлечения}

\subsubsection{Функциональность}
Компонент поиска и извлечения обрабатывает пользовательские запросы и возвращает релевантные сегменты видео.

\subsubsection{Технические Детали}
\begin{itemize}
    \item \textbf{Обработка Запросов}
    \begin{itemize}
        \item Реализация нормализации и расширения запросов
        \item Преобразование естественного языка в структурированные запросы
        \item Поддержка запросов взаимосвязей концепций
    \end{itemize}

    \item \textbf{Поиск по Индексу}
    \begin{itemize}
        \item Оптимизация операций булева поиска
        \item Реализация близостного поиска для запросов с несколькими концепциями
        \item Применение фильтрации по релевантности
    \end{itemize}

    \item \textbf{Ранжирование Результатов}
    \begin{itemize}
        \item Применение сложного алгоритма ранжирования
        \item Извлечение точных сегментов видео
        \item Генерация фрагментов результатов
    \end{itemize}
\end{itemize}

\subsubsection{Конвейер Обработки Запросов}
\begin{figure}[H]
\centering
\begin{tikzpicture}[
    node distance=0.3cm,
    auto,
    block/.style={rectangle, draw, fill=blue!10, text width=3.7cm, text centered, rounded corners, minimum height=0.7cm},
    line/.style={draw, -stealth, thick}
]
    % First row
    \node [block] (input) {Пользовательский запрос};
    \node [block, right=of input] (norm) {Нормализация запроса};
    \node [block, right=of norm] (expand) {Расширение запроса};

    % Second row
    \node [block, below=of input, yshift=-0.2cm] (struct) {Создание структурированного запроса};
    \node [block, right=of struct] (index) {Поиск по индексу};
    \node [block, right=of index] (rank) {Ранжирование результатов};

    % Third row
    \node [block, below=of struct, yshift=-0.2cm] (context) {Извлечение контекста};
    \node [block, right=of context, text width=8cm] (output) {Ранжированный список сегментов видео с временными метками};

    % Connections
    \path [line] (input) -- (norm);
    \path [line] (norm) -- (expand);
    \path [line] (expand) |- ([yshift=-0.3cm]expand.south) -| (struct);
    \path [line] (struct) -- (index);
    \path [line] (index) -- (rank);
    \path [line] (rank) |- ([yshift=-0.3cm]rank.south) -| (context);
    \path [line] (context) -- (output);
\end{tikzpicture}
\caption{Конвейер обработки запросов}
\end{figure}

\subsection{Компонент Интеграции}

\subsubsection{Функциональность}
Компонент интеграции предоставляет API для подключения к существующей платформе.

\subsubsection{Технические Детали}
\begin{itemize}
    \item \textbf{Конечные Точки REST API}
    \begin{itemize}
        \item Реализация конечной точки отправки видео
        \item Создание интерфейса поисковых запросов
        \item Предоставление возможностей мониторинга статуса
    \end{itemize}

    \item \textbf{Асинхронная Обработка}
    \begin{itemize}
        \item Реализация управления очередью обработки
        \item Предоставление обновлений статуса в реальном времени
        \item Обработка восстановления после сбоев
    \end{itemize}

    \item \textbf{Интеграция с Фронтендом}
    \begin{itemize}
        \item Определение структурированного формата ответов
        \item Поддержка расширенных параметров запроса
        \item Реализация стратегии кэширования на стороне клиента
    \end{itemize}
\end{itemize}

\subsubsection{Спецификации API}
\begin{lstlisting}[language=json, caption=Формат ответа поискового API]
{
  "query": "Weierstrass theorem",
  "total_results": 24,
  "results": [
    {
      "video_id": "abcd1234",
      "title": "Advanced Calculus - Lecture 5",
      "concept": "Weierstrass theorem",
      "relevance_score": 89,
      "coverage_type": "substantive",
      "timestamp_start": 1250.5,
      "timestamp_end": 1380.2,
      "context_snippet": "Today we're going to prove the Weierstrass theorem which states that...",
      "thumbnail_url": "https://...",
      "related_concepts": ["uniform convergence", "continuous functions"]
    },
    // More results...
  ],
}
\end{lstlisting}

\section{План Реализации}

\begin{table}[H]
\centering
\begin{tabular}{@{}llp{7cm}@{}}
\toprule
\textbf{Фаза} & \textbf{Продолжительность} & \textbf{Результаты} \\
\midrule
1. Базовая Инфраструктура & неделя 1 &
  \begin{itemize}[leftmargin=*, nosep]
    \item Интеграция с YouTube API
    \item Конвейер извлечения транскриптов
    \item Реализация алгоритма MLCS на C++
  \end{itemize} \\
\midrule
2. Анализ Контента & неделя 2 &
  \begin{itemize}[leftmargin=*, nosep]
    \item Конвейер предобработки NLP
    \item Система извлечения концепций
    \item Реализация контекстного анализа
  \end{itemize} \\
\midrule
3. Система Индексации & недели 3-4 &
  \begin{itemize}[leftmargin=*, nosep]
    \item Реализация инвертированного индекса
    \item Алгоритмы оценки релевантности
    \item Слой сохранения индекса
  \end{itemize} \\
\midrule
4. Поисковая Система & недели 5-6 &
  \begin{itemize}[leftmargin=*, nosep]
    \item Реализация обработки запросов
    \item Алгоритм ранжирования результатов
    \item Конечные точки поискового API
    \item Реализация демонстративного консольного интерфейса
  \end{itemize} \\
\midrule
5. Интеграция & в будущем &
  \begin{itemize}[leftmargin=*, nosep]
    \item Интеграция с бэкендом Spring
    \item Компоненты фронтенда React
    \item Система асинхронной обработки
  \end{itemize} \\
\midrule
6. Тестирование \& Оптимизация & в будущем &
  \begin{itemize}[leftmargin=*, nosep]
    \item Тестирование производительности
    \item Настройка релевантности
    \item Исправление ошибок и оптимизация
  \end{itemize} \\
\bottomrule
\end{tabular}
\caption{Фазы разработки}
\end{table}

\section{Технические Требования}

\subsection{Требования к Производительности}

\begin{itemize}
    \item \textbf{Эффективность Обработки}
    \begin{itemize}
        \item Обработка часовой видеолекции < 5 минут
        \item Поддержка пакетной обработки нескольких видео
        \item Максимальное использование ОЗУ < 4 ГБ
    \end{itemize}

    \item \textbf{Производительность Поиска}
    \begin{itemize}
        \item Время отклика запроса < 200 мс для 95\% запросов
        \item Поддержка не менее 50 одновременных пользователей
        \item Размер индекса растет сублинейно с объемом контента
    \end{itemize}

    \item \textbf{Качество Релевантности}
    \begin{itemize}
        \item Точность > 0.8 для классификации содержательного освещения
        \item Полнота > 0.9 для идентификации концепций
    \end{itemize}
\end{itemize}

\section{Заключение}

Разработанный ''Индексатор содержания видеолекций'' позволит студентам быстро находить нужные концепции в массиве видеоматериалов, четко различая случайные упоминания и полноценные объяснения тем.

В основе системы лежит применение алгоритма Linear MLCS в сочетании с методами обработки естественного языка. Модульная структура проекта обеспечивает его последовательную реализацию с четкими промежуточными результатами. Использование C++ для высокопроизводительных компонентов и Python для обработки данных создает оптимальный баланс между быстродействием и гибкостью разработки.

По завершении проект интегрируется с параллельно разрабатываемой студенческой платформой, существенно расширяя её функциональность за счет интеллектуального поиска по видеоконтенту.

\appendix
\section{Глоссарий}

\begin{itemize}
    \item \textbf{Linear MLCS (Множественная Наибольшая Общая Подпоследовательность)} - Алгоритм для нахождения наибольшей общей подпоследовательности в нескольких последовательностях.

    \item \textbf{Инвертированный Индекс} - Структура данных, отображающая контент (концепции) на их расположение (видео, временные метки).

    \item \textbf{Глубина Освещения} - Мера того, насколько тщательно объяснена или обсуждена концепция.

    \item \textbf{Стабильность Концепции} - То, как долго концепция остается в центре внимания обсуждения в видео.

    \item \textbf{Дискурсивные Маркеры} - Лингвистические сигналы, указывающие на объяснение, определение или ссылку.

    \item \textbf{Распознавание Сущностей (NER)} - Техника NLP для идентификации именованных сущностей в тексте.
\end{itemize}

\end{document}
